% === Begin Label Data ===
% ID: 792
% 难度星级: 
% 标签: 
% 备注: 
% 组卷引用次数: 0
% === End  Label Data ===

\begin{problem}{2016}{G}{全国I卷}{7}{三角函数}
已知函数 $ f(x)=\sin(\omega x+\varphi) $ ($ \omega>0 $, $ |\varphi|\leqslant\dfrac{\pi}{2} $), $ x=-\dfrac{\pi}{4} $ 为 $ f(x) $ 的零点, $ x=\dfrac{\pi}{4} $ 为函数 $ y=f(x) $ 图像的对称轴, 且 $ f(x) $ 在 $ \left(\dfrac{\pi}{18},\dfrac{5\pi}{36}\right) $ 单调, 则 $ \omega $ 的最大值为 (\hspace{1cm})
\begin{choices}
\choice{{$11$}}
\choice{{$9$}}
\choice{{$7$}}
\choice{{$5$}}
\end{choices}
\end{problem}

\begin{answer}
B
\end{answer}

\begin{solutions}
思路1：由正弦型三角函数的单调性推出 $ \omega $ 满足的关系.  

因为 $ f(x) $ 在 $ \left(\dfrac{\pi}{18},\dfrac{5\pi}{36}\right) $ 单调, 所以区间 $ \left(\dfrac{\pi}{18},\dfrac{5\pi}{36}\right) $ 不能包含函数 $ f(x) $ 的最值点, 即  
$$
\dfrac{5\pi}{36}-\dfrac{\pi}{18}\leqslant\dfrac{\pi}{\omega},\quad\text{化简得}\omega\leqslant12.
$$  

因为 $ x=-\dfrac{\pi}{4} $ 为 $ f(x) $ 的零点, $ x=\dfrac{\pi}{4} $ 为 $ y=f(x) $ 图像的对称轴, 所以  
$$
\begin{cases}
-\dfrac{\pi}{4}\omega+\varphi=m\pi\,(m\in\mathbb{Z}),\\[1em]
\dfrac{\pi}{4}\omega+\varphi=n\pi+\dfrac{\pi}{2}\,(m\in\mathbb{Z}),
\end{cases}
$$  
由此可得 $ \varphi=\dfrac{k\pi}{2}+\dfrac{\pi}{4}\,(k\in\mathbb{Z}) $.  

又 $ |\varphi|\leqslant\dfrac{\pi}{2} $, 故 $ \omega=\dfrac{\pi}{4} $ 或 $ \varphi=-\dfrac{\pi}{4} $.  

当 $ \varphi=\dfrac{\pi}{4} $ 时, $ \omega=1-4m\,(m\in\mathbb{Z}) $, 而 $ \omega>0 $ 且 $ \omega\leqslant12 $, 可得 $ \omega $ 的可能值为 $ 1,5,9 $;  

当 $ \varphi=-\dfrac{\pi}{4} $ 时, $ \omega=-1-4m\,(m\in\mathbb{Z}) $, 而 $ \omega>0 $ 且 $ \omega\leqslant12 $, 可得 $ \omega $ 的可能值为 $ 3,7,11 $;  

验证 $ f(x)=\sin\left(11x-\dfrac{\pi}{4}\right) $ 有一个最值点 $ x=\dfrac{3\pi}{44}\in\left(\dfrac{\pi}{18},\dfrac{5\pi}{36}\right) $, 不满足题设.  

验证 $ f(x)=\sin\left(9x+\dfrac{\pi}{4}\right) $ 满足题设, 故选 B.  

思路2：由正弦型三角函数的零点及对称轴分析 $ \omega $ 与 $ \varphi $ 满足的条件.  

因为 $ x=-\dfrac{\pi}{4} $ 为 $ f(x) $ 的零点, $ x=\dfrac{\pi}{4} $ 为 $ y=f(x) $ 图像的对称轴, 所以  
$$
\begin{cases}
-\dfrac{\pi}{4}\omega+\varphi=m\pi\,(m\in\mathbb{Z}),\\[1em]
\dfrac{\pi}{4}\omega+\varphi=n\pi+\dfrac{\pi}{2}\,(m\in\mathbb{Z}),
\end{cases}
$$  
由此可得 $ \omega=1-4m\,(m\in\mathbb{Z}) $, $ \omega=1+4n\,(n\in\mathbb{Z}) $, $ \varphi=\dfrac{t\pi}{2}+\dfrac{\pi}{4}\,(t\in\mathbb{Z}) $.  

又 $ \omega>0 $, $ |\varphi|\leqslant\dfrac{\pi}{2} $, 故 $ \omega=\dfrac{\pi}{4} $ 或 $ \varphi=-\dfrac{\pi}{4} $.  

因为 $ f(x) $ 在 $ \left(\dfrac{\pi}{18},\dfrac{5\pi}{36}\right) $ 单调, 所以区间 $ \left(\dfrac{\pi}{18},\dfrac{5\pi}{36}\right) $ 不能包含函数 $ f(x) $ 的最值点, 即  
$$
\dfrac{5\pi}{36}-\dfrac{\pi}{18}\leqslant\dfrac{\pi}{\omega},\quad\text{化简得}\omega\leqslant12.
$$  

验证 $ f(x)=\sin\left(11x-\dfrac{\pi}{4}\right) $ 有一个最值点 $ x=\dfrac{3\pi}{44}\in\left(\dfrac{\pi}{18},\dfrac{5\pi}{36}\right) $, $ f(x)=\sin\left(7x-\dfrac{\pi}{4}\right) $ 有一个最值点 $ x=\dfrac{3\pi}{28}\in\left(\dfrac{\pi}{18},\dfrac{5\pi}{36}\right) $, $ f(x)=\sin\left(3x-\dfrac{\pi}{4}\right) $ 满足题设.  

验证 $ f(x)=\sin\left(9x+\dfrac{\pi}{4}\right) $, $ f(x)=\sin\left(5x+\dfrac{\pi}{4}\right) $, $ f(x)=\sin\left(x+\dfrac{\pi}{4}\right) $ 满足题设.
\end{solutions}